\section{Background}
People show emotion through several channels at once: the face, the voice, and the electrical activity of the brain. Affective computing sets out to read these signals automatically, with uses ranging from human-computer interaction and education to mental-health monitoring \cite{Picard1997AffectiveComputing}. It has usually been framed as classification---from the available cues, decide which emotion is present.

That framing quietly assumes that the emotion a person shows is the emotion they feel. When someone is asked to act out an emotion in a laboratory, the assumption is reasonable. Outside it, people routinely hide or tone down what they feel: a patient may appear composed while distressed, and people with alexithymia or depression may display little of what they experience. Signals that are harder to control deliberately, such as EEG, might therefore add information that the face and voice do not carry.

These observations motivate two questions that run through this thesis. The first is a fusion question: when three modalities of very different character---a voluntary facial expression, a voluntary vocal expression, and an involuntary neural signal---are combined, does a learned, attention-based fusion module extract more from them than simply averaging their predictions, and under what conditions? The second is a coherence question: when the expressed signals (audio and video) and the neural signal (EEG) disagree, is that disagreement structured in a way that reflects the person rather than the classifiers?

The experimental vehicle for both is the EAV dataset \cite{Lee2024EAV}, a laboratory corpus of 42 participants with 30-channel EEG, audio and video recorded across five emotion categories: Neutral, Anger, Happiness, Sadness and Calmness. Three trimodal fusion methods have been published on it, all evaluated within-subject.

\section{Motivation}
The motivation for this research emerges from three observations: one architectural, one methodological, and one clinical.

The architectural observation is that the modalities in a corpus like EAV are not equivalent. Audio and video are external signals under substantial voluntary control; EEG is an internal signal that is largely outside it. They also differ in reliability: a facial-emotion encoder pre-trained on large image corpora starts far ahead of an EEG classifier trained from scratch on a few hundred trials. A fusion module that weights modalities dynamically---per input, as attention does---should in principle exploit these differences better than fixed averaging. Whether it does is an empirical question, and the answer on EAV has not been measured carefully.

The methodological observation is that how a multimodal emotion system is evaluated can matter as much as how it is built. Evaluating on the same participants a system was trained on measures something different from evaluating on new ones, and choosing a model on the data used to score it inflates the score. Both issues are known in principle; their size on EAV was unknown.

The clinical observation is that a mismatch between shown and felt emotion has been linked in the literature to emotion suppression and with conditions of impaired emotional expressivity such as alexithymia and depression \cite{Gross1998EmotionRegulation,Mauss2009Measures}. A tool that could detect such discordance from passively recorded multimodal data would be clinically relevant---provided the discordance it reports is a property of the person and not of its own classifiers. The present thesis is scoped to a healthy laboratory population and makes no clinical claim.

\section{Problem Statement}
Four inter-related problems define the challenge this thesis addresses.

First, the \textbf{fusion problem}. Learned fusion modules are routinely reported to outperform simple late fusion, but those comparisons are only as good as the protocol behind them. On a corpus with 280 training trials per participant, a 2.29-million-parameter attention module may have too little data to learn more than averaging provides, and a fair comparison requires that neither side benefits from selection on the test set.

Second, the \textbf{generalisation problem}. Every published trimodal result on EAV is within-subject: the model has seen the test participant's face, voice and neural signature during training. How much of that accuracy survives on a new person, and which modality loses it, is unknown. If the loss is caused by participant-specific structure in the features, the question becomes whether it can be removed without labels.

Third, the \textbf{signal-quality problem in EEG}. EEG recorded during natural emotional stimuli is a low signal-to-noise measurement, contaminated by muscle and motion artefacts and non-stationary across sessions and participants. With 280 training trials per participant and no cross-subject pretraining, EEG classifier confidence is informative only when it is high; the coherence analysis therefore requires an EEG top-1 confidence of at least 0.5.

Fourth, the \textbf{corpus-realism problem}. EAV is cue-based: participants are told which emotion to express. This enables reliable labels, but it means that disagreement between modalities may reflect imperfect simulation rather than suppression, that the audio-visual consensus is usually the cued label, and that the EEG (recorded while listening) and the audio and video (recorded while speaking) are not simultaneous. Any coherence claim must survive these confounds.

\section{Objectives and Research Questions}
The thesis addresses nine research questions. RQ1--RQ5 were posed in Pre-Thesis~II; RQ6 and RQ7 were added when the subject-independent evaluation was undertaken; RQ8 and RQ9 were added in the final phase.

\begin{enumerate}
    \item \textbf{RQ1---Per-modality baseline.} How well does each modality alone classify the five EAV emotions, compared with the dataset paper's baselines (audio 36.7\%, vision 52.8\%, EEG 36.7\%)?
    \item \textbf{RQ2---Fusion gain.} By how much, if at all, does combining the modalities beat the strongest one on its own?
    \item \textbf{RQ3---Architectural contribution.} Does learned trimodal attention fusion outperform a naive softmax-averaging late-fusion baseline when neither is selected on the test set?
    \item \textbf{RQ4---Coherence as signal.} When audio-visual consensus and EEG disagree, what patterns emerge, and do they exceed what the classifiers' own independent errors would produce?
    \item \textbf{RQ5---Robustness.} Does the system produce useful predictions when EEG is unavailable at inference time?
    \item \textbf{RQ6---Cross-subject generalisation.} Does the system recognise emotions in participants it has never seen during training, and how much does accuracy change relative to the within-subject protocol?
    \item \textbf{RQ7---Protocol dependence of the architectural contribution.} Does the comparison between learned fusion and naive averaging survive subject-independent evaluation?
    \item \textbf{RQ8---Subject calibration.} Can a label-free, per-participant normalisation of the fusion input, computed from a short unlabelled recording, restore the advantage of learned fusion on unseen participants; how much calibration data does it need; and how sensitive is it to the emotional content of that recording?
    \item \textbf{RQ9---Evaluation integrity.} Which conclusions of the Pre-Thesis~II analysis survive leak-free model selection, a stronger trial-alignment check, and a permutation test of the coherence patterns?
\end{enumerate}

\section{Contributions}
\label{sec:contributions_intro}
The thesis makes two principal contributions, supported by three secondary ones.

\begin{enumerate}
    \item \textbf{Subject calibration for trimodal fusion (Chapter~\ref{chap:calibration}).} Standardising each participant's fusion features with statistics from about 20 unlabelled five-second clips lifts every learned fusion head on unseen participants from 3--4.5 points below averaging to 73.6--75.9\%. Linear probes show that the calibration removes linearly decodable participant identity, mostly from vision. Averaging calibrated the same way is just as accurate, so without labels the gain belongs to calibration rather than to the fusion operator; when the calibration clips are also labelled, the learned heads outperform labelled averaging by 3.2--4.8 points at 50 clips, with attention trained with modality dropout the most accurate.
    \item \textbf{An evaluation-integrity audit (Chapter~\ref{chap:audit}).} Choosing each fusion head's epoch on the test set had inflated the Pre-Thesis~II fusion results by 5.3--7.5 points; the modalities' trials cannot be shown to correspond from the released data; and the suppression matrix cannot be told apart from the EEG classifier's own errors. Each test is quick to run and can be reused on other work with EAV.
    \item \textbf{The first subject-independent evaluation of trimodal fusion on EAV (Chapter~\ref{chap:crosssubject})}, which traces the loss on new participants to the vision modality.
    \item \textbf{A trimodal attention fusion module that works without EEG (Chapters~\ref{chap:methodology} and~\ref{chap:within})}, trained with softhard modality dropout so that accuracy holds when EEG is absent.
    \item \textbf{Documented fixes to the EAV reference code (Chapter~\ref{chap:engineering_challenges})}, including defects in the public EEGNet code that prevented it from training correctly.
\end{enumerate}

\section{Methodology in Brief}
The system is built in three cached stages, and evaluated under four protocols.

\textbf{Stage A: per-modality encoders.} The Audio Spectrogram Transformer (AST), pre-trained on AudioSet, is fine-tuned on audio clips; a Vision Transformer (ViT) pre-trained for facial emotion is fine-tuned on video frames; and EEGNet is trained from scratch on EEG segments.

\textbf{Stage B: feature extraction.} Each trained encoder's pre-classifier representation is extracted for every trial, giving one pooled feature vector per modality per trial (768, 768 and 960 dimensions).

\textbf{Stage C: fusion.} Four fusion methods are compared on the same features: naive late fusion (the mean of the three modalities' softmax outputs); a trimodal attention module (\texttt{TrimodalAttentionFusion}) that treats the three modality features as a three-token sequence and applies self-attention; a concatenation MLP; and the attention module trained with softhard modality dropout, which is also evaluated with EEG zeroed.

\textbf{Protocols.} Results are reported within-subject (each participant's 280/120 split), cross-subject (five-fold subject-wise cross-validation with encoders retrained per fold), cross-subject with per-participant calibration, and cross-subject with labelled calibration clips. Every trained model's reported epoch is chosen without reference to the data it is scored on. Paired Wilcoxon tests with Holm correction and bootstrap confidence intervals are used throughout.

\textbf{Coherence analysis.} A suppression matrix counts trials on which audio and vision agree, EEG confidently disagrees, and records the pair of emotions involved. Its cells are tested against a class-conditional permutation null that preserves every classifier's error rates.

\section{Scope and Challenges}
\textbf{Scope.} The thesis is scoped to clip-level classification of five-second segments into the five EAV emotion categories, on all 42 EAV participants. No new data were recorded: a planned EEG acquisition was cancelled for budget reasons, and all EEG comes from EAV. The intended deployment configuration uses audio and video only, with EEG zeroed, matching any realistic setting without an EEG headset; the zero-EEG results of Chapters~\ref{chap:within} and~\ref{chap:calibration} evaluate that configuration.

\textbf{Challenges.} Twenty-four technical and methodological challenges are documented in Chapter~\ref{chap:engineering_challenges}. The most consequential were defects in the EAV reference EEGNet code that prevented it from training correctly; the leakage risk of reusing within-subject encoders for a cross-subject evaluation, which required retraining every encoder for every fold; and the discovery, in the final phase, that the Pre-Thesis~II fusion results had been selected on the test set.

\section{Report Organisation}
Chapter~\ref{chap:litreview} reviews the relevant literature. Chapter~\ref{chap:requirements} sets out requirements, impacts and constraints. Chapter~\ref{chap:methodology} describes the methodology and its implementation. Chapter~\ref{chap:within} reports within-subject results; Chapter~\ref{chap:crosssubject}, the subject-independent evaluation; Chapter~\ref{chap:calibration}, unsupervised subject calibration; and Chapter~\ref{chap:audit}, the evaluation-integrity audit. Chapter~\ref{chap:engineering_challenges} documents engineering challenges, and Chapter~\ref{chap:conclusion} concludes.
